% Chapter 4 -- Diskettes and the Fixed Disk.
% Sources (Law 1): os/milton/command.c builtin_vol; os/milton/dosutil.c
% label_main; os/milton/chkdsk.c; os/milton/xcopy_tree.c (XCOPY into a new
% directory, names shown, keeps date and time); spec/hardware.json storage
% (FAT12 diskette image, FAT16 fixed-disk image, ATA PIO).  FDISK, FORMAT,
% DISKCOPY and DISKCOMP are shipped as DOS\FDISK.COM, DOS\FORMAT.COM,
% DOS\DISKCOPY.COM and DOS\DISKCOMP.COM (Makefile FLAIR_DOS_EXTERNALS;
% os/milton/fdisk.c, os/milton/diskutil.c; WL-0141, WL-0147) and have their
% entries in the Reference.  DISKCOMP after COPY: MS-DOS 3.3 User's Guide,
% Reference pp.59-62 (MS-DOS_3.3_Users_Guide_198707.pdf PDF 144-147),
% "Compare OK" and the first Note (initech-090x).
% GATES: test-label-vol test-dosutil-emu test-chkdsk test-chkdsk-emu test-xcopy-tree-emu test-fat test-pubs-screens
\Gates{test-label-vol, test-dosutil-emu, test-chkdsk, test-chkdsk-emu, test-xcopy-tree-emu, test-fat, test-pubs-screens}
\chapter{Diskettes and the Fixed Disk}

\section{Kinds of Disk}
\index{diskettes}\index{fixed disk}
InitechDOS keeps files on two kinds of disk.  A \term{diskette} is a
flexible disk in a protective jacket that you put in a diskette drive and
take out again.  A \term{fixed disk} is built into the computer\Slip{VO30-07}{, it} holds
far more than a diskette and is much faster.  InitechDOS reads and writes
both in the same way, and the commands in this book work on either.

In this version of InitechDOS the disk that holds your files is drive~A,
whichever kind of disk it is.

\section{Taking Care of Diskettes}
\begin{itemize}
\item Hold a diskette by its label.  Do not touch the recording surface.
\item Keep diskettes away from magnets, telephones and the display, and
  out of direct sunlight.
\item Write on a diskette's label before you put it on the diskette, or
  use a felt-tip pen.
\item Do not remove a diskette from its drive while the drive light is on.
\item Store diskettes upright in their envelopes, in a box.
\end{itemize}

\section{Preparing a New Disk}
\index{formatting}\index{DISKCOPY@\texttt{DISKCOPY}}\index{DISKCOMP@\texttt{DISKCOMP}}
A new diskette must be \term{formatted} before InitechDOS can use it:
formatting marks out the disk and creates an empty root directory.  A new
fixed disk must first be divided into partitions with \cmd{FDISK}, then
formatted.  The \cmd{FORMAT} command does both kinds of formatting, and
\cmd{DISKCOPY} makes an exact copy of a whole diskette.  Formatting erases
everything on a disk.

To check a copy, use \cmd{DISKCOMP}.  It compares the two diskettes sector
by sector and displays \texttt{Compare OK} when they are the same.  With one
diskette drive, both commands ask for each diskette in turn.  A diskette
whose files were copied one by one with \cmd{COPY} may not compare, even
though the files are the same: \cmd{COPY} does not put the files in the same
places on the disk.  The \doctitle{InitechDOS Reference} describes the
options of all four commands.

\section{Volume Labels}
\index{volume label}\index{VOL@\texttt{VOL}}\index{LABEL@\texttt{LABEL}}
A volume label names a disk, so that you can tell your disks apart.  The
\cmd{VOL} command displays the label:
\begin{Steps}
\item Type \type{CLS} and press \key{Enter}.
\item Type \type{VOL} and press \key{Enter}.
\end{Steps}

\Screen[rows=4, label=fig:vol]{screens/vol.txt}{Displaying the volume label}

The \cmd{LABEL} utility program changes the label.  Typed alone, it shows
the current label and asks for a new one, of up to eleven characters:
\begin{Steps}
\item Type \type{CLS} and press \key{Enter}.
\item Type \type{LABEL} and press \key{Enter}.
\item At the question, type \type{WORKDISK} and press \key{Enter}.
\item Type \type{VOL} and press \key{Enter}.
\end{Steps}

\Screen[rows=7, label=fig:label]{screens/label.txt}{Changing the volume label}

You can also give the label on the command line: \type{LABEL INITECH} puts
the original label back.  If you press \key{Enter} without typing a label,
\cmd{LABEL} asks whether to delete the current label.

\section{Checking a Disk}
\index{checking a disk}\index{CHKDSK@\texttt{CHKDSK}}
The \cmd{CHKDSK} utility program checks that the directories and the file
allocation table of a disk agree, and reports how the disk and the memory of
the computer are used.  Run it from time to time, and whenever a program
has stopped while it was writing to the disk.
\begin{Steps}
\item Type \type{CLS} and press \key{Enter}.
\item Type \type{CHKDSK} and press \key{Enter}.
\end{Steps}

\Screen[rows=12, label=fig:chkdsk]{screens/chkdsk.txt}{A disk with no errors}

The report gives the space on the disk, how much of it is used by hidden
files, directories and your own files, how much is still available, and the
total and free memory.  When \cmd{CHKDSK} finds errors it describes them.
Run it again as \type{CHKDSK /F} to correct them; Chapter~9 and the
\doctitle{InitechDOS Reference} describe what it does.

\section{Making Backup Copies}
\index{backup copies}\index{XCOPY@\texttt{XCOPY}}
A disk can be damaged or lost, and a file can be deleted by mistake.  Keep
copies of your important files on another disk, and bring them up to date
regularly.  The \cmd{XCOPY} utility program copies groups of files, and
whole directories, in one command.  In this exercise you copy the text files
into a new directory, \file{BACKUP}:
\begin{Steps}
\item Type \type{CLS} and press \key{Enter}.
\item Type \type{XCOPY *.TXT BACKUP\bs} and press \key{Enter}.  The
  backslash at the end tells \cmd{XCOPY} that \file{BACKUP} is a directory;
  it creates it.
\item Type \type{CD BACKUP} and press \key{Enter}, then \type{DIR}.
\end{Steps}

\Screen[rows=20, label=fig:backup]{screens/backup.txt}{Copying the text files with XCOPY}

\cmd{XCOPY} lists each file as it copies it.  Each copy keeps the date and
time of its original, so the listing shows when each file was last
changed.  To copy only the files that have changed since your last backup,
use \cmd{XCOPY} with \kw{/M} (Chapter~9).  The new directory
\file{BACKUP} and its \file{.} and \file{..} entries carry the date and time
at which \cmd{XCOPY} made it.
